% ==============================================================================
% Fichier  : 05_theorie_information_shannon.tex
% Titre    : Fondements de la Théorie de l'Information de Shannon
%            appliqués à l'Audit des Systèmes d'Information
% Cours    : Audit des Systèmes d'Information – ENSP, Département GI
% Niveau   : Humanités Numériques, Niveau 3
% ==============================================================================

\chapter{Fondements de la Théorie de l'Information de Shannon}
\label{chap:shannon}

% ──────────────────────────────────────────────────────────────────────────────
\section{Pourquoi un auditeur doit-il connaître la théorie de l'information~?}
\label{sec:motivation}
% ──────────────────────────────────────────────────────────────────────────────

L'audit des systèmes d'information est, dans son essence, une activité de
réduction d'incertitude. L'auditeur arrive dans une organisation dont il ignore
l'état réel~: les contrôles sont-ils en place~? La documentation existe-t-elle~?
Les processus sont-ils appliqués~? Chaque question posée, chaque document
examiné, chaque entretien conduit contribue à réduire cette incertitude initiale
et à construire une opinion argumentée sur la qualité du contrôle interne.

Cette réduction progressive de l'incertitude n'est pas seulement une métaphore
intuitive~: elle peut être \emph{mesurée}. C'est précisément l'objet de la
théorie mathématique de l'information, fondée par Claude Elwood Shannon en
1948~\cite{Shannon1948}. Cette théorie fournit une métrique rigoureuse, l'entropie,
qui quantifie l'incertitude contenue dans une distribution de probabilités. Dans
le chapitre suivant, cette métrique sera directement mise en œuvre pour
quantifier l'incertitude résiduelle d'un audit à travers la distribution des
réponses trichotomiques O~/~N~/~PAS. Comprendre son fondement théorique est donc
une condition préalable à l'utilisation correcte et à l'interprétation juste de
l'outil d'aide à l'audit.

Ce chapitre présente les notions fondamentales dans l'ordre de leur construction
logique~: la quantification de l'information (§\,\ref{sec:information}), l'entropie
de Shannon et ses propriétés (§\,\ref{sec:entropie}), la notion d'information
mutuelle et de réduction d'incertitude (§\,\ref{sec:info-mutuelle}), puis
l'application directe à la mesure de l'incertitude dans un audit
(§\,\ref{sec:application-audit}).

\begin{boxlizbiz}
\textbf{Cas Lizbiz's~: l'incertitude comme objet d'audit}

Lors d'une mission préliminaire chez Lizbiz's, l'auditeur senior dresse un bilan
intermédiaire~: sur~30~questions du domaine Sécurité, 10~ont reçu une réponse~O,
10~une réponse~N, et 10~une réponse~PAS. L'équipe dispose donc d'une vision à peu
près aussi informative que le hasard le plus complet. Ce constat, avant même
d'avoir produit un rapport, est \emph{en soi} un constat d'audit~: une
organisation dont les acteurs ne savent pas si ses propres contrôles de sécurité
fonctionnent est une organisation en situation de risque structurel. La mesure
de cette situation par l'entropie de Shannon permettra de la rendre communicable,
comparable et traçable d'un exercice à l'autre.
\end{boxlizbiz}

% ──────────────────────────────────────────────────────────────────────────────
\section{La mesure de l'information}
\label{sec:information}
% ──────────────────────────────────────────────────────────────────────────────

\subsection{Le problème de la surprise}

La question fondatrice de Shannon est d'une simplicité trompeuse~: \emph{combien
d'information apporte un message~?} L'intuition quotidienne suggère que la réponse
dépend du degré de surprise provoqué par ce message. Apprendre que le soleil s'est
levé ce matin n'apporte aucune information~: c'est un événement certain. Apprendre
qu'un système d'information national a subi une intrusion réussie malgré une
politique de sécurité certifiée ISO~27001 est fortement informatif~: c'est un
événement de faible probabilité.

Cette intuition peut être formalisée~: l'information apportée par un événement
doit être une fonction \emph{décroissante} de sa probabilité. Shannon montre
qu'à deux contraintes supplémentaires près, continuité de la fonction et
additivité pour les événements indépendants, il existe une unique famille de
fonctions satisfaisant ces propriétés~\cite{Shannon1948, Cover2006}.

\begin{boxdef}{Information propre (auto-information)}
Soit $\mathcal{X}$ un espace de probabilité fini et $x \in \mathcal{X}$ un
événement de probabilité $p(x) \in (0,\,1]$. L'\emph{information propre}
(ou \emph{auto-information}) de $x$ est définie par~:
\begin{equation}
  I(x) \;=\; -\log_2 p(x) \;\in\; [0,\,+\infty)
  \label{eq:autopinfo}
\end{equation}
L'unité de mesure est le \emph{bit} lorsque le logarithme est en base~2.
\end{boxdef}

Quelques cas limites illustrent la cohérence de cette définition~:
\begin{itemize}
  \item si $p(x) = 1$ (événement certain), alors $I(x) = -\log_2 1 = 0$~: un
        événement certain n'apporte aucune information~;
  \item si $p(x) = 1/2$, alors $I(x) = 1$~bit~: l'issue d'un tirage à pile-ou-face
        équitable apporte exactement~1~bit d'information~;
  \item si $p(x) \to 0$ (événement rarissime), alors $I(x) \to +\infty$~: un événement
        de probabilité quasi nulle serait, s'il se réalisait, porteur d'une
        information quasi infinie.
\end{itemize}

\paragraph{Propriété d'additivité.}
Si $x$ et $y$ sont deux événements \emph{indépendants}, l'information apportée
par leur réalisation conjointe est la somme des informations individuelles~:
\begin{equation}
  I(x,y) \;=\; I(x) + I(y)
  \label{eq:additivite}
\end{equation}
Cette propriété est fondamentale~: elle est la traduction formelle du fait qu'apprendre
deux nouvelles indépendantes est aussi informatif qu'apprendre chacune séparément.

\subsection{De l'événement à la source}

Dans la pratique, on s'intéresse rarement à un événement isolé mais à une
\emph{source d'information}, c'est-à-dire une variable aléatoire $X$ prenant ses
valeurs dans un ensemble fini $\mathcal{X} = \{x_1, \ldots, x_n\}$ avec des
probabilités $p(x_i) = \Pr[X = x_i]$. La question pertinente devient alors~:
quelle est l'incertitude \emph{moyenne} sur la valeur que va prendre $X$, avant
qu'on en connaisse la réalisation~? La réponse est l'entropie de Shannon.

% ──────────────────────────────────────────────────────────────────────────────
\section{L'entropie de Shannon}
\label{sec:entropie}
% ──────────────────────────────────────────────────────────────────────────────

\subsection{Définition}

\begin{boxdef}{Entropie de Shannon}
Soit $X$ une variable aléatoire discrète à valeurs dans $\mathcal{X} =
\{x_1, \ldots, x_n\}$, de loi $p$. L'\emph{entropie de Shannon} de $X$ est la
valeur d'information propre \emph{moyenne} sous la loi $p$~:
\begin{equation}
  H(X) \;=\; \mathbb{E}_{p}\bigl[I(X)\bigr]
           \;=\; -\sum_{i=1}^{n} p(x_i)\,\log_2 p(x_i)
  \label{eq:entropie}
\end{equation}
avec la convention $0 \log_2 0 = 0$ (limite par continuité). L'unité est le
\emph{bit}. $H(X)$ mesure l'incertitude moyenne associée à la source $X$.
\end{boxdef}

La formulation de Shannon est remarquable par sa concision~: elle réduit la notion
d'incertitude, intuitive mais diffuse, à une seule quantité scalaire, calculable,
additive et reliée à la notion de compression optimale de l'information
\cite{Shannon1948, MacKay2003}.

\subsection{Propriétés fondamentales}

Les quatre propriétés suivantes sont des théorèmes démontrables à partir de la
définition~\eqref{eq:entropie}~\cite{Cover2006}.

\begin{boiteTheoreme}{Non-négativité et borne supérieure de l'entropie}
Pour toute variable aléatoire discrète $X$ à $n$ valeurs~:
\begin{equation}
  0 \;\leq\; H(X) \;\leq\; \log_2 n
  \label{eq:bornes}
\end{equation}
\begin{itemize}
  \item $H(X) = 0$ si et seulement si $X$ est \emph{déterministe}~: une et une
        seule valeur $x_i$ a une probabilité~1 (certitude complète, incertitude
        nulle)~;
  \item $H(X) = \log_2 n$ si et seulement si $X$ est \emph{uniformément distribuée}~:
        toutes les valeurs sont équiprobables (incertitude maximale).
\end{itemize}
\end{boiteTheoreme}

\paragraph{Application au contexte trichotomique de l'audit.}
Dans l'outil d'aide à l'audit (chapitre~\ref{chap:06_outil_aide_audit}), les réponses
prennent trois valeurs~: $n = 3$. La borne supérieure de l'entropie est donc
$\log_2 3 \approx 1{,}585$~bits, valeur atteinte lorsque les trois modalités O, N
et PAS/NEK sont équiprobables. C'est le repère cardinal de l'outil~: une entropie
proche de~$1{,}585$ signale que les réponses d'un point de contrôle sont aussi
incertaines que le hasard le plus parfait; une entropie proche de~0 signale une
situation de quasi-certitude (toutes les réponses convergent vers la même
modalité).

\begin{boiteTheoreme}{Concavité de l'entropie}
La fonction $H$ est \emph{strictement concave} sur le simplexe des distributions
de probabilité~: pour toutes distributions $p$ et $q$ et tout $\lambda \in (0,1)$,
\begin{equation}
  H\bigl(\lambda p + (1-\lambda) q\bigr) \;\geq\; \lambda H(p) + (1-\lambda) H(q)
  \label{eq:concavite}
\end{equation}
\end{boiteTheoreme}

La concavité a une interprétation directe~: mélanger deux sources d'information
ne peut qu'augmenter ou conserver l'incertitude, jamais la réduire. En termes
d'audit, agréger les réponses de deux équipes d'auditeurs indépendants sur le même
périmètre ne peut que maintenir ou accroître l'entropie globale~; elle ne saurait
la réduire.

\begin{boiteTheoreme}{Sous-additivité}
Pour deux variables aléatoires $X$ et $Y$, l'entropie jointe satisfait~:
\begin{equation}
  H(X, Y) \;\leq\; H(X) + H(Y)
  \label{eq:sous-additivite}
\end{equation}
avec égalité si et seulement si $X$ et $Y$ sont \emph{indépendantes}.
\end{boiteTheoreme}

\begin{boiteTheoreme}{Entropie conditionnelle et règle de chaînage}
L'\emph{entropie conditionnelle} de $X$ sachant $Y$ est l'incertitude résiduelle
sur $X$ lorsque $Y$ est connue~:
\begin{equation}
  H(X \mid Y) \;=\; -\sum_{j} p(y_j) \sum_{i} p(x_i \mid y_j)\,\log_2 p(x_i \mid y_j)
  \label{eq:entropie-cond}
\end{equation}
La \emph{règle de chaînage} s'écrit alors~:
\begin{equation}
  H(X, Y) \;=\; H(X) + H(Y \mid X) \;=\; H(Y) + H(X \mid Y)
  \label{eq:chainage}
\end{equation}
On a toujours $H(X \mid Y) \leq H(X)$~: la connaissance d'une observation $Y$
ne peut qu'égaler ou réduire l'incertitude sur $X$, jamais l'augmenter.
\end{boiteTheoreme}

\subsection{Illustration sur des distributions élémentaires}

Le tableau~\ref{tab:entropies} calcule l'entropie pour plusieurs distributions
trichotomiques représentatives de situations d'audit.

\begin{table}[H]
\centering
\small
\caption{Entropie de Shannon pour des distributions trichotomiques {O, N, PAS}
représentatives de situations d'audit}
\label{tab:entropies}
\begin{tabular}{lccccc}
\toprule
\textbf{Situation d'audit} & $p_O$ & $p_N$ & $p_P$ &
$H$ (bits) & \textbf{Interprétation} \\
\midrule
Conformité totale &
1{,}00 & 0{,}00 & 0{,}00 & 0{,}000 & Certitude~: aucune zone d'ombre \\
\addlinespace
Forte non-conformité &
0{,}10 & 0{,}80 & 0{,}10 & 0{,}922 & Situation dégradée mais lisible \\
\addlinespace
Forte incertitude &
0{,}10 & 0{,}10 & 0{,}80 & 0{,}922 & Carence informationnelle majeure \\
\addlinespace
Incertitude modérée &
0{,}50 & 0{,}25 & 0{,}25 & 1{,}500 & Zone grise significative \\
\addlinespace
Incertitude maximale &
$\frac{1}{3}$ & $\frac{1}{3}$ & $\frac{1}{3}$ & 1{,}585 & Chaos informationnel \\
\bottomrule
\end{tabular}
\end{table}

Deux observations méritent d'être soulignées. Premièrement, une forte
non-conformité avérée (distribution $0{,}10$-$0{,}80$-$0{,}10$) et une forte
carence documentaire (distribution $0{,}10$-$0{,}10$-$0{,}80$) produisent la
\emph{même} entropie~: les deux situations sont, du point de vue informationnel,
également certaines, bien qu'elles appellent des recommandations distinctes.
Deuxièmement, la situation de zone grise modérée ($0{,}50$-$0{,}25$-$0{,}25$),
qui pourrait paraître rassurante du fait de la majorité de réponses O, est
presque aussi entropique que le désordre complet~: la présence de~$25\,\%$ de
réponses PAS/NEK suffit à élever fortement l'incertitude globale.

\begin{boiteRemarque}{Calcul de l'entropie~: règle de convention}
Par définition, $0 \log_2 0 = 0$. Cette convention est indispensable pour que
la formule~\eqref{eq:entropie} reste valide lorsqu'une ou plusieurs modalités
ont une probabilité nulle. Dans un tableur, elle se traduit par l'usage d'un
test conditionnel~: \verb|SI(p=0;0;p*LOG(p;2))|. Omettre cette convention produit
une erreur de division par zéro.
\end{boiteRemarque}

% ──────────────────────────────────────────────────────────────────────────────
\section{Information mutuelle et réduction d'incertitude}
\label{sec:info-mutuelle}
% ──────────────────────────────────────────────────────────────────────────────

\subsection{Définition}

La notion d'information mutuelle formalise la question~: \emph{combien une
observation $Y$ réduit-elle l'incertitude sur $X$~?} C'est la traduction
quantitative de l'activité même de l'audit.

\begin{boxdef}{Information mutuelle}
L'\emph{information mutuelle} de deux variables aléatoires $X$ et $Y$ est
définie par~:
\begin{equation}
  I(X\,;\,Y) \;=\; H(X) - H(X \mid Y) \;=\; H(Y) - H(Y \mid X)
  \label{eq:info-mutuelle}
\end{equation}
Elle mesure la réduction d'incertitude sur $X$ apportée par la connaissance de
$Y$, et réciproquement. On a $I(X\,;\,Y) \geq 0$, avec égalité si et seulement
si $X$ et $Y$ sont indépendantes.
\end{boxdef}

L'information mutuelle s'exprime également comme la divergence de
Kullback-Leibler~\cite{Kullback1951} entre la loi jointe $p(x,y)$ et le produit
des lois marginales $p(x)\,p(y)$~:
\begin{equation}
  I(X\,;\,Y) \;=\; \sum_{x,y} p(x,y)\,\log_2 \frac{p(x,y)}{p(x)\,p(y)}
  \label{eq:kl}
\end{equation}

\subsection{Interprétation dans le contexte de l'audit}

La table~\ref{tab:analogie-audit} établit le parallèle entre les notions de la
théorie de l'information et les réalités de la mission d'audit.

\begin{table}[H]
\centering
\small
\caption{Correspondance entre théorie de l'information et pratique de l'audit}
\label{tab:analogie-audit}
\begin{tabular}{p{4.5cm}p{4.5cm}p{4.5cm}}
\toprule
\textbf{Notion formelle} & \textbf{Notation} & \textbf{Réalité de l'audit} \\
\midrule
Variable aléatoire $X$ &
$X$ : état réel du contrôle &
Le contrôle est-il effectif~? L'auditeur l'ignore avant l'audit. \\
\addlinespace
Distribution a priori de $X$ &
$p(x)$ avant investigation &
Croyances initiales de l'auditeur~; en l'absence de données~: 
distribution uniforme sur O/N/PAS. \\
\addlinespace
Observation $Y$ &
Réponse obtenue (O, N, PAS) &
Réponse à une question d'évaluation, document examiné, entretien conduit. \\
\addlinespace
Entropie a priori $H(X)$ &
Incertitude avant investigation &
Incertitude totale à l'ouverture de la mission. \\
\addlinespace
Entropie a posteriori $H(X \mid Y)$ &
Incertitude après observation &
Incertitude résiduelle après collecte des preuves. \\
\addlinespace
Information mutuelle $I(X\,;\,Y)$ &
$H(X) - H(X \mid Y)$ &
Valeur informative d'une question d'audit~: réduction d'incertitude produite
par l'investigation. \\
\bottomrule
\end{tabular}
\end{table}

La mission d'audit peut donc se lire comme un processus de maximisation de
l'information mutuelle~: l'auditeur cherche, par son plan d'investigation, à
maximiser $I(X\,;\,Y_1, \ldots, Y_k)$, c'est-à-dire à sélectionner les questions
et les observations qui réduisent le plus l'incertitude initiale $H(X)$. Un
point de contrôle pour lequel l'entropie reste élevée après investigation, parce
que toutes les réponses sont PAS/NEK, est un point pour lequel l'information
mutuelle entre les observations et l'état réel du contrôle est nulle~: les
questions posées n'ont apporté aucune réduction d'incertitude. C'est un constat
d'audit en soi.

\subsection{La divergence de Kullback-Leibler et la mesure de l'écart normatif}

La théorie de l'information offre un outil complémentaire~: la divergence de
Kullback-Leibler (KL), qui mesure l'écart entre une distribution observée et une
distribution de référence~\cite{Kullback1951}.

\begin{boxdef}{Divergence de Kullback-Leibler}
Soient $p$ (distribution observée) et $q$ (distribution de référence) deux
distributions sur un espace fini $\mathcal{X}$. La \emph{divergence KL} de $p$
par rapport à $q$ est~:
\begin{equation}
  D_{\mathrm{KL}}(p \,\|\, q) \;=\; \sum_{x \in \mathcal{X}} p(x)\,\log_2\frac{p(x)}{q(x)}
  \label{eq:kl-div}
\end{equation}
avec $D_{\mathrm{KL}}(p \,\|\, q) \geq 0$ et égalité si et seulement si $p = q$
(inégalité de Gibbs~\cite{Cover2006}).
\end{boxdef}

Dans le contexte de l'audit, $q$ peut représenter la distribution cible d'une
organisation mature, par exemple $q = (0{,}90,\,0{,}05,\,0{,}05)$ pour
un domaine bien contrôlé, et $p$ la distribution observée lors d'une mission.
La divergence KL mesure alors l'écart entre la réalité observée et l'état
souhaité~: c'est un indicateur synthétique de la distance entre la maturité
effective et la maturité cible.

\begin{boiteRemarque}{Non-symétrie de la divergence KL}
La divergence KL n'est pas symétrique~: $D_{\mathrm{KL}}(p \,\|\, q) \neq
D_{\mathrm{KL}}(q \,\|\, p)$ en général. Elle ne définit donc pas une distance
au sens mathématique strict. Pour des usages où la symétrie est souhaitable
(comparaison entre deux unités auditées), on lui préfère la divergence de Jensen-Shannon,
définie comme la moyenne des divergences KL de $p$ et $q$ vers leur mélange
$m = \frac{1}{2}(p+q)$~\cite{Lin1991}. Son usage reste optionnel dans le cadre de
ce cours~; la divergence KL suffit pour les applications présentées au chapitre
suivant.
\end{boiteRemarque}

% ──────────────────────────────────────────────────────────────────────────────
\section{Le codage de l'information~: lien avec la compression}
\label{sec:codage}
% ──────────────────────────────────────────────────────────────────────────────

L'entropie possède une interprétation opérationnelle fondamentale qui en justifie
le nom et en éclaire le sens profond~: elle représente la \emph{longueur minimale
moyenne} d'un code nécessaire pour représenter les messages d'une source. Ce
résultat est le \emph{théorème du codage source} de Shannon.

\begin{boiteTheoreme}{Théorème du codage source (Shannon, 1948)}
Soit $X$ une source discrète sans mémoire de loi $p$. Tout code préfixe
(code de Huffman~\cite{Huffman1952}, par exemple) attribue à chaque symbole $x_i$
un mot de code de longueur $\ell_i$. La \emph{longueur moyenne} $\bar{L} =
\sum_i p(x_i)\,\ell_i$ satisfait~:
\begin{equation}
  H(X) \;\leq\; \bar{L} \;<\; H(X) + 1
  \label{eq:codage-source}
\end{equation}
La borne inférieure $H(X)$ est atteinte asymptotiquement~: l'entropie est la
limite fondamentale de la compression sans perte.
\end{boiteTheoreme}

L'interprétation pour l'audit est la suivante. Si un point de contrôle présente
une entropie $H \approx 0$, ses réponses sont quasi certaines~: elles sont peu
informatives dans l'absolu, mais elles sont \emph{compressibles} en très peu de
bits. À l'inverse, une entropie $H \approx 1{,}585$~bits signifie que les
réponses de ce point de contrôle sont si imprévisibles qu'elles nécessitent la
quasi-totalité du potentiel informationnel d'un système à trois symboles~: chaque
réponse est véritablement une \emph{nouvelle} au sens informationnel du terme.

Cette intuition rejoint directement la pratique de l'audit~: un point de contrôle
à forte entropie mérite plus de temps, d'investigations complémentaires et de
diligences documentaires, précisément parce que chaque réponse y est plus
susceptible d'être surprenante et donc potentiellement révélatrice.

% ──────────────────────────────────────────────────────────────────────────────
\section{Entropie et distributions trichotomiques~: analyse exhaustive}
\label{sec:trichotomique}
% ──────────────────────────────────────────────────────────────────────────────

\subsection{Paramétrisation de la surface d'entropie}

La surface d'entropie d'une distribution trichotomique $(p_O, p_N, p_P)$ avec
$p_O + p_N + p_P = 1$ forme un simplexe bidimensionnel. Son étude complète
permet de comprendre comment chacun des trois types de réponse contribue à l'entropie
globale.

En posant $\mathbf{p} = (p_O, p_N, p_P)$, on a~:
\begin{equation}
  H(\mathbf{p})
    = -p_O \log_2 p_O - p_N \log_2 p_N - p_P \log_2 p_P
    \;\in\; [0,\; \log_2 3]
  \label{eq:entropie-tri}
\end{equation}

Les lignes de niveau de cette surface révèlent trois zones de qualité distinctes~:

\begin{enumerate}
  \item \textbf{Zone de certitude} ($H < 0{,}5$~bit)~: au moins l'une des trois
        modalités est largement dominante. L'audit est conclusif, même si la
        conclusion est défavorable. Un audit entièrement rouge (toutes les
        réponses~N) est, paradoxalement, aussi informatif qu'un audit entièrement
        vert~: dans les deux cas, $H = 0$.
  \item \textbf{Zone d'incertitude modérée} ($0{,}5 \leq H < 1{,}0$~bit)~:
        les réponses sont partagées mais avec une tendance perceptible.
        L'investigation complémentaire est utile mais non urgente.
  \item \textbf{Zone grise critique} ($H \geq 1{,}0$~bit)~: les trois modalités
        coexistent de façon significative. L'audit de ce point de contrôle doit
        être étendu avant que toute conclusion soit émise.
\end{enumerate}

\subsection{Décomposition de l'entropie et contribution des composantes}

L'entropie d'une distribution trichotomique peut se décomposer en la somme de
l'entropie binaire O~vs.~(N+PAS) et de l'entropie conditionnelle sur la
répartition N/PAS~:
\begin{equation}
  H(p_O, p_N, p_P) \;=\;
    H_b(p_O) \;+\; (1 - p_O)\,H\!\left(\frac{p_N}{1-p_O},\;
    \frac{p_P}{1-p_O}\right)
  \label{eq:decomp}
\end{equation}
où $H_b(\alpha) = -\alpha \log_2 \alpha - (1-\alpha)\log_2(1-\alpha)$ est
l'entropie binaire de Bernoulli. Cette décomposition est utile pour identifier
la source principale d'incertitude~: le premier terme mesure l'incertitude sur la
dichotomie conformité/non-conformité, tandis que le second mesure l'incertitude
entre non-conformité avérée et incertitude documentaire.

% ──────────────────────────────────────────────────────────────────────────────
\section{Application à la mesure de l'incertitude dans un audit}
\label{sec:application-audit}
% ──────────────────────────────────────────────────────────────────────────────

\subsection{Modélisation formelle du processus d'audit}

Soit $\mathcal{D}$ un domaine d'audit décomposé en $K$ points de contrôle
$\{C_1, \ldots, C_K\}$, et soit $C_k$ un point de contrôle comportant $n_k$
questions d'évaluation. Pour chaque question $j \in \{1, \ldots, n_k\}$, l'auditeur
collecte une réponse $r_{k,j} \in \{\texttt{O},\,\texttt{N},\,\texttt{PAS}\}$.

\begin{boxdef}{Entropie d'un point de contrôle}
Notons $n_k^O$, $n_k^N$, $n_k^P$ les effectifs de chaque modalité dans le point
de contrôle $C_k$, et $n_k = n_k^O + n_k^N + n_k^P$ le nombre total de questions.
L'\emph{entropie empirique} du point de contrôle est~:
\begin{equation}
  H(C_k) \;=\; -\!\!\!\sum_{m \,\in\, \{\texttt{O},\texttt{N},\texttt{P}\}}
  \hat{p}_{k,m}\,\log_2 \hat{p}_{k,m}
  \label{eq:entropie-cp}
\end{equation}
où $\hat{p}_{k,m} = n_k^m / n_k$ est la fréquence observée de la modalité~$m$.
\end{boxdef}

\begin{boxdef}{Entropie d'un domaine (entropie agrégée)}
L'\emph{entropie du domaine} $\mathcal{D}$ est la moyenne pondérée des entropies
de ses points de contrôle, pondérée par le nombre de questions~:
\begin{equation}
  H(\mathcal{D}) \;=\; \frac{\displaystyle\sum_{k=1}^{K} n_k\,H(C_k)}{\displaystyle\sum_{k=1}^{K} n_k}
  \label{eq:entropie-domaine}
\end{equation}
Cette quantité est l'entropie empirique calculée sur l'ensemble des réponses du
domaine, traitées comme un seul échantillon.
\end{boxdef}

\subsection{Le double signal de l'outil d'audit}

Le chapitre~\ref{chap:06_outil_aide_audit} introduit l'idée que l'outil mesure deux
dimensions distinctes~: \emph{ce qui ne va pas} (taux de non-conformité,
proportion de réponses~N) et \emph{ce qu'on ne sait pas} (proportion de réponses
PAS/NEK). La théorie de Shannon éclaire pourquoi ces deux dimensions sont
fondamentalement irréductibles l'une à l'autre.

Formellement, les deux signaux sont~:
\begin{align}
  \tau_c(\mathcal{D}) &= \hat{p}_O \;=\; \frac{n^O}{n}
  \label{eq:taux-conformite} \\[6pt]
  H(\mathcal{D}) &= -\hat{p}_O \log_2 \hat{p}_O
                   - \hat{p}_N \log_2 \hat{p}_N
                   - \hat{p}_P \log_2 \hat{p}_P
  \label{eq:entropie-domaine2}
\end{align}

Ces deux indicateurs sont \emph{orthogonaux}~: deux domaines de même taux de
conformité peuvent présenter des entropies très différentes, et inversement. Le
tableau~\ref{tab:ortho} illustre cette orthogonalité.

\begin{table}[H]
\centering
\small
\caption{Orthogonalité du taux de conformité et de l'entropie~: quatre
configurations distinctes pour un même taux de conformité de~$50\,\%$}
\label{tab:ortho}
\begin{tabular}{lccccc}
\toprule
\textbf{Configuration} & $\hat{p}_O$ & $\hat{p}_N$ & $\hat{p}_P$ &
$\tau_c$ & $H$ (bits) \\
\midrule
Audit clair, moitié conforme & 0{,}50 & 0{,}50 & 0{,}00 & 50\,\% & 1{,}000 \\
Audit ambigu, sans zone grise & 0{,}50 & 0{,}25 & 0{,}25 & 50\,\% & 1{,}500 \\
Forte carence documentaire & 0{,}50 & 0{,}10 & 0{,}40 & 50\,\% & 1{,}361 \\
Incertitude maximale & $\frac{1}{3}$ & $\frac{1}{3}$ & $\frac{1}{3}$ & 33\,\% & 1{,}585 \\
\bottomrule
\end{tabular}
\end{table}

Ces quatre situations n'ont pas le même profil de risque, bien que les deux
premières partagent le même taux de conformité de~$50\,\%$. Dans la première,
l'auditeur sait ce qui va et ce qui ne va pas~; dans la seconde, un quart des
réponses reste dans la zone grise, ce qui signifie qu'un quart du périmètre
audité est soustrait à l'opinion de l'auditeur. Réduire l'audit à son seul taux
de conformité reviendrait à traiter ces situations de façon identique, ce qui
serait une erreur méthodologique grave.

\subsection{Interprétation normative des seuils d'entropie}

La table~\ref{tab:seuils} résume les seuils d'interprétation de l'entropie en
contexte d'audit, tels qu'ils seront mis en œuvre dans l'outil du chapitre
suivant.

\begin{table}[H]
\centering
\small
\caption{Seuils normatifs d'entropie pour l'interprétation d'un point de
contrôle dans l'outil d'aide à l'audit} 
\label{tab:seuils}
\begin{tabular}{cccl}
\toprule
\textbf{Plage} & \textbf{Couleur outil} & \textbf{Qualification} &
\textbf{Conduite à tenir} \\
\midrule
$H < 0{,}5$ & Blanc & Certitude & Aucune investigation complémentaire requise \\
$0{,}5 \leq H < 1{,}0$ & Jaune & Incertitude modérée &
Relance de l'audité sur les PAS/NEK \\
$H \geq 1{,}0$ & Rouge & Zone grise critique &
Investigation approfondie obligatoire avant conclusion \\
\bottomrule
\end{tabular}
\end{table}

\begin{boiteRemarque}{Nature statistique de l'entropie empirique}
L'entropie~$H(C_k)$ calculée dans l'outil est une \emph{entropie empirique}~:
elle est estimée sur un échantillon fini de $n_k$ questions. Pour des points de
contrôle comportant peu de questions ($n_k < 10$), l'estimation peut être biaisée
et instable. L'estimateur du maximum de vraisemblance de l'entropie est
asymptotiquement non biaisé mais surestimé pour de petits $n$~\cite{Paninski2003}.
Des correcteurs de biais (correcteur de Miller-Madow) existent mais ne sont pas
nécessaires dans le cadre de ce cours, où l'entropie est utilisée comme indicateur
de pilotage et non comme estimateur statistique précis.
\end{boiteRemarque}

% ──────────────────────────────────────────────────────────────────────────────
\section{Extensions~: entropie croisée et mesures apparentées}
\label{sec:extensions}
% ──────────────────────────────────────────────────────────────────────────────

Pour les lecteurs souhaitant aller plus loin, deux extensions de la notion
d'entropie de Shannon sont présentées~: l'entropie croisée, utilisée pour mesurer
l'adéquation d'un modèle à une distribution observée, et l'entropie de Rényi,
qui généralise l'entropie de Shannon à une famille paramétrée.

\begin{boxdef}{Entropie croisée}
Soient $p$ (distribution réelle) et $q$ (distribution modèle) deux distributions
sur $\mathcal{X}$. L'\emph{entropie croisée} de $p$ relativement à $q$ est~:
\begin{equation}
  H(p, q) \;=\; -\sum_{x \in \mathcal{X}} p(x)\,\log_2 q(x)
  \label{eq:entropie-croisee}
\end{equation}
On a $H(p,q) = H(p) + D_{\mathrm{KL}}(p\,\|\,q) \geq H(p)$, avec égalité si
et seulement si $p = q$.
\end{boxdef}

En audit, $q$ peut représenter la distribution cible d'une organisation certifiée
(ex.~$q_{\mathrm{cible}} = (0{,}90, 0{,}05, 0{,}05)$) et $p$ la distribution observée.
L'entropie croisée $H(p, q_{\mathrm{cible}})$ mesure alors le coût informationnel
de l'inadéquation entre la réalité et la cible~: plus ce coût est élevé, plus
l'organisation s'écarte de son niveau de maturité attendu.

\begin{boxdef}{Entropie de Rényi}
Soit $\alpha > 0$, $\alpha \neq 1$. L'\emph{entropie de Rényi d'ordre $\alpha$}
est~:
\begin{equation}
  H_\alpha(X) \;=\; \frac{1}{1-\alpha}\,\log_2\!\left(\sum_{i} p(x_i)^\alpha\right)
  \label{eq:renyi}
\end{equation}
À la limite $\alpha \to 1$, $H_\alpha(X) \to H(X)$ (entropie de Shannon).
\end{boxdef}

L'entropie de Rényi d'ordre~$\alpha > 1$ pénalise davantage les événements
fréquents et est moins sensible aux événements rares~; pour $\alpha < 1$,
l'effet inverse se produit. Dans le contexte de l'audit, l'entropie de Rényi
d'ordre~2 ($H_2$), appelée \emph{entropie de collision}, offre une mesure
d'incertitude plus robuste aux petites erreurs d'estimation des probabilités~;
elle peut être utilisée comme vérification de la stabilité des résultats obtenus
avec l'entropie de Shannon standard~\cite{Renyi1961}.

% ──────────────────────────────────────────────────────────────────────────────
\section{Synthèse~: de Shannon à l'audit}
\label{sec:synthese}
% ──────────────────────────────────────────────────────────────────────────────

Ce chapitre a construit le socle théorique nécessaire à la pleine compréhension
de l'outil d'aide à l'audit présenté au chapitre suivant. La progression logique
en est la suivante~:

\begin{enumerate}
  \item L'information propre $I(x) = -\log_2 p(x)$ quantifie la surprise associée
        à l'observation d'un événement isolé~; son analogue en audit est l'apport
        informationnel d'une seule réponse d'évaluation.
  \item L'entropie de Shannon $H(X)$ est la moyenne de cette surprise sur
        l'ensemble d'une source~; appliquée à un point de contrôle, elle mesure
        l'incertitude résiduelle de l'auditeur après investigation.
  \item L'information mutuelle $I(X\,;\,Y)$ mesure la réduction d'incertitude
        apportée par une observation~; elle formalise la valeur informative d'un
        programme d'audit~: un bon plan d'audit est celui qui maximise la
        réduction d'incertitude pour un coût d'investigation donné.
  \item Le théorème du codage source relie l'entropie à la compressibilité~;
        un point de contrôle à entropie élevée est, au sens formel, un point
        dont les réponses sont fondamentalement incompressibles~: elles ne peuvent
        être prédites ni remplacées par des connaissances préalables.
  \item La divergence de Kullback-Leibler permet de mesurer l'écart entre l'état
        observé d'un domaine et un référentiel de maturité cible~; elle est
        l'instrument formel de la notation de maturité.
\end{enumerate}

\begin{boiteTheoreme}{Principe de la cartographie entropique en audit}
Soit un audit portant sur $K$ points de contrôle. Définissons le vecteur
d'entropies $\mathbf{H} = (H(C_1), \ldots, H(C_K)) \in [0,\,\log_2 3]^K$.
La \emph{cartographie entropique} de l'audit est la représentation graphique de
$\mathbf{H}$. Elle constitue un second tableau de bord orthogonal au vecteur des
taux de conformité $\boldsymbol{\tau} = (\tau_{c,1}, \ldots, \tau_{c,K})$.

Les quatre quadrants du plan $(\tau_c, H)$ correspondent à quatre pathologies
distinctes~:
\begin{itemize}
  \item $(\tau_c$ élevé, $H$ faible)~: bonne conformité, haute lisibilité~---
        la situation nominale visée~;
  \item $(\tau_c$ faible, $H$ faible)~: non-conformité massive mais clairement
        identifiée, situation grave mais traçable~;
  \item $(\tau_c$ élevé, $H$ élevée)~: apparence de conformité masquée par une
        carence d'information, situation trompeuse, à risque sous-estimé~;
  \item $(\tau_c$ faible, $H$ élevée)~: non-conformité probable et incertitude
        élevée, la situation la plus préoccupante, requérant une investigation
        immédiate.
\end{itemize}
\end{boiteTheoreme}

\begin{boxlizbiz}
\textbf{Cas Lizbiz's~: lecture de la cartographie entropique}

À l'issue de l'audit de Lizbiz's, la feuille \texttt{SYNTHÈSE} produit le
vecteur d'entropies suivant pour cinq domaines~:

\smallskip
\begin{tabular}{lccl}
\toprule
\textbf{Domaine} & $\tau_c$ & $H$ (bits) & \textbf{Quadrant} \\
\midrule
Applications & $78\,\%$ & $0{,}62$ & Bonne conformité, lisibilité moyenne \\
DSI & $55\,\%$ & $0{,}38$ & Non-conformité marquée, lisible \\
Projets & $82\,\%$ & $1{,}21$ & \textbf{Conformité apparente, fort risque masqué} \\
Support & $91\,\%$ & $0{,}18$ & Situation nominale \\
Sécurité & $33\,\%$ & $1{,}40$ & \textbf{Non-conformité et incertitude critiques} \\
\bottomrule
\end{tabular}

\smallskip
La lecture du seul taux de conformité aurait conduit à concentrer les
recommandations sur les domaines DSI et Sécurité. La cartographie entropique
révèle que le domaine Projets, apparemment sain ($\tau_c = 82\,\%$), est en
réalité une zone grise critique ($H = 1{,}21$~bits)~: plus de la moitié des
réponses de ce domaine sont PAS/NEK, signalant une organisation de projet dont
personne ne connaît le fonctionnement réel. Ce domaine exige une investigation
complémentaire prioritaire avant que l'auditeur puisse émettre une opinion.
\end{boxlizbiz}

% ──────────────────────────────────────────────────────────────────────────────
\section*{Références bibliographiques}
% ──────────────────────────────────────────────────────────────────────────────

\addcontentsline{toc}{section}{Références bibliographiques}

\begin{itemize}

\item \textsc{Shannon}, C.\,E. (1948). A mathematical theory of communication.
\emph{Bell System Technical Journal}, 27(3), 379--423 \& 27(4), 623--656.
\href{https://doi.org/10.1002/j.1538-7305.1948.tb01338.x}{DOI}

\item \textsc{Cover}, T.\,M. et \textsc{Thomas}, J.\,A. (2006).
\emph{Elements of Information Theory} (2\textsuperscript{e}~éd.).
Wiley-Interscience. ISBN 978-0-471-24195-9.

\item \textsc{MacKay}, D.\,J.\,C. (2003).
\emph{Information Theory, Inference, and Learning Algorithms}.
Cambridge University Press. ISBN 978-0-521-64298-9.
Disponible en libre accès~: \url{http://www.inference.org.uk/mackay/itila/}

\item \textsc{Kullback}, S. et \textsc{Leibler}, R.\,A. (1951).
On information and sufficiency.
\emph{The Annals of Mathematical Statistics}, 22(1), 79--86.
\href{https://doi.org/10.1214/aoms/1177729694}{DOI: 10.1214/aoms/1177729694}

\item \textsc{Huffman}, D.\,A. (1952).
A method for the construction of minimum-redundancy codes.
\emph{Proceedings of the IRE}, 40(9), 1098--1101.
\href{https://doi.org/10.1109/JRPROC.1952.273898}{DOI: 10.1109/JRPROC.1952.273898}
\item \textsc{Lin}, J. (1991).
Divergence measures based on the Shannon entropy.
\emph{IEEE Transactions on Information Theory}, 37(1), 145--151.
\href{https://doi.org/10.1109/18.61115}{DOI: 10.1109/18.61115}

\item \textsc{Rényi}, A. (1961).
On measures of entropy and information.
\emph{Proceedings of the 4th Berkeley Symposium on Mathematical Statistics and
Probability}, 1, 547--561.
University of California Press.

\item \textsc{Paninski}, L. (2003).
Estimation of entropy and mutual information.
\emph{Neural Computation}, 15(6), 1191--1253.
\href{https://doi.org/10.1162/089976603321780272}{DOI: 10.1162/089976603321780272}

\item \textsc{Jaynes}, E.\,T. (1957).
Information theory and statistical mechanics.
\emph{Physical Review}, 106(4), 620--630.
\href{https://doi.org/10.1103/PhysRev.106.620}{DOI: 10.1103/PhysRev.106.620}

\item \textsc{ISACA} (2022).
\emph{CISA Review Manual, 28\textsuperscript{th} Edition}.
ISACA. ISBN 978-1-60420-981-4.

\end{itemize}

% ──────────────────────────────────────────────────────────────────────────────
\section*{Exercices}
% ──────────────────────────────────────────────────────────────────────────────

\addcontentsline{toc}{section}{Exercices}

\begin{boxexo}
\textbf{Exercice~1 – Calcul direct d'entropie}

Un auditeur clôture l'examen d'un point de contrôle \emph{Gestion des accès}
comportant 12~questions. Les réponses sont~: 6~O, 3~N, 3~PAS.

\begin{enumerate}[label=(\alph*)]
  \item Calculer l'entropie empirique $H$ du point de contrôle.
  \item Situer cette valeur dans la table de seuils (table~\ref{tab:seuils}).
  \item Conclure sur la conduite à tenir.
\end{enumerate}
\end{boxexo}

\begin{boxexo}
\textbf{Exercice~2 – Orthogonalité des indicateurs}

Deux domaines $\mathcal{D}_1$ et $\mathcal{D}_2$ présentent le même taux de
conformité $\tau_c = 60\,\%$. La distribution de $\mathcal{D}_1$ est
$(0{,}60, 0{,}30, 0{,}10)$ et celle de $\mathcal{D}_2$ est $(0{,}60, 0{,}10, 0{,}30)$.

\begin{enumerate}[label=(\alph*)]
  \item Calculer $H(\mathcal{D}_1)$ et $H(\mathcal{D}_2)$.
  \item Ces deux valeurs sont-elles différentes~? Justifier.
  \item Quel domaine présente le risque le plus préoccupant du point de vue de
        la qualité de l'information~? Argumenter.
\end{enumerate}
\end{boxexo}

\begin{boxexo}
\textbf{Exercice~3 – Divergence de Kullback-Leibler et écart normatif}

Une organisation de référence présente la distribution cible
$q = (0{,}90, 0{,}05, 0{,}05)$. L'organisation Lizbiz's présente sur le domaine
Sécurité la distribution observée $p = (0{,}33, 0{,}40, 0{,}27)$.

\begin{enumerate}[label=(\alph*)]
  \item Calculer $D_{\mathrm{KL}}(p \,\|\, q)$.
  \item Interpréter le résultat en termes d'écart de maturité.
  \item Commenter la non-symétrie en calculant $D_{\mathrm{KL}}(q \,\|\, p)$.
\end{enumerate}
\end{boxexo}

\begin{boxexo}
\textbf{Exercice~4 – Codage et interprétation}

Un point de contrôle présente une distribution de réponses $(0{,}80, 0{,}15, 0{,}05)$.

\begin{enumerate}[label=(\alph*)]
  \item Calculer l'entropie $H$.
  \item En utilisant le théorème du codage source, donner la longueur minimale
        moyenne d'un code binaire représentant les réponses de ce point de
        contrôle.
  \item Un code naïf assignerait~2~bits par symbole (codage en base~4). Quel
        est le gain de compression théorique possible~?
\end{enumerate}
\end{boxexo}

